% !TeX root = ../guide
\chapter{TeXstudio: A Comprehensive IDE for \LaTeX}\label{ch:texstudio}
	\section{Introduction}\label{sec:TSIntroduction}
		While \LaTeX\ documents can be written in any plain text editor, using a dedicated Integrated Development Environment (IDE) like TeXstudio significantly lowers the learning curve. 
		TeXstudio is designed specifically for creating \LaTeX\ documents, offering features like real-time syntax checking, integrated PDF viewing, and automated command completion.
		This chapter will provide an overview of TeXstudio's primary features and guide you through setting it up to work seamlessly with your thesis template.
	
	\section{Key Features of TeXstudio}
		TeXstudio stands out because it brings together all the tools needed for document preparation into a single interface, reducing the need to interact directly with the command line.
		
		\subsection{Interactive Syntax Checking}
			One of the most useful features for students is the \enquote{live} syntax checker. 
			TeXstudio underlines errors in red and warnings in yellow as you type. 
			This allows you to catch a missing bracket or a misspelled command immediately, rather than waiting until you try to compile the entire document.
		
		\subsection{Integrated PDF Viewer}
			TeXstudio includes a built-in PDF viewer that supports \enquote{SyncTeX.} 
			This allows for two-way communication between your code and your output: if you right-click a sentence in the PDF, you can jump directly to the corresponding line in the \LaTeX\ code, and vice versa.
		
		\subsection{Advanced Completion}
			When you begin typing a command (\textit{e.g.}, \texttt{\textbackslash cite\{}), TeXstudio provides a dropdown list of available citation keys from your bibliography. 
			It also automatically suggests the required arguments for complex environments like tables and figures, which is a massive time-saver for a thesis.
		
		\subsection{Structure View}
			The side panel provides a hierarchical view of your document structure (Parts, Chapters, Sections). 
			For a large thesis, this serves as an interactive Table of Contents, allowing you to jump between different chapters of your work with a single click.
	
	\section{Getting Started with TeXstudio}
		Setting up TeXstudio correctly is essential for ensuring that your thesis compiles without errors, especially when using the uAlberta template.
		
		\subsection{Installation}
			The software is available for Windows, macOS, and Linux. 
			Visit the official website: \url{https://www.texstudio.org/} and download the installer for your operating system.
			\note{%
				Ensure you have already installed a \LaTeX\ distribution like MikTeX or TeXLive before installing TeXstudio, as the editor requires a compiler to function.}
		
		\subsection{The Main Interface}
			Upon launching TeXstudio, you will see the primary workspace as shown in \Cref{fig:TeXstudioMain}.
			\begin{figure}[htbp]
				\centering
				\includegraphics[width=0.8\linewidth]{example-image}
				\caption{TeXstudio Main Interface showing editor, structure, and PDF viewer.}
				\label{fig:TeXstudioMain}
			\end{figure}
		
		\subsection{Configuration for Your Thesis}
			To ensure the uAlberta template compiles correctly (especially if you are using specialized fonts or Bib\LaTeX), you should verify your build settings:
			\begin{enumerate}
				\item Go to `Options' $\rightarrow$ `Configure TeXstudio'.
				\item Under the `Build' tab, ensure the `Default Compiler' is set to \textbf{PdfLaTeX}.
				\item Ensure the `Default Bibliography Tool' is set to \textbf{Biber}.
			\end{enumerate}

	\section{Efficient Writing Workflows}
		Beyond just typing text, TeXstudio has several tools that make the technical side of thesis writing easier.
		
		\subsection{Using the Bibliography with JabRef}
			If you have your JabRef library open and saved as a \path{.bib} file in your project folder, TeXstudio will automatically \enquote{read} it. 
			When you type \texttt{\textbackslash cite\{}, a list of your references will appear. 
			You can search this list by author or title directly within the editor.
		
		\subsection{Managing Large Documents}
			Since a thesis is usually split into multiple files (one for each chapter), you must tell TeXstudio which file is the \enquote{Root} file.
			This can be done in one of two ways.

			The first is to use \enquote{Magic Comments} at the top of each separate file.
			The format of this comment (which should be the first line of the file) should be:
			\texttt{\% !TeX root = ../ualberta.tex}.

			The second way is to set the root in TeXstudio.
			\begin{enumerate}
				\item Open your main thesis file (\textit{e.g.}, \path{ualberta.tex}).
				\item Go to `Options' $\rightarrow$ `Root Document' $\rightarrow$ `Set Current Document as Explicit Root'.
			\end{enumerate}
			This ensures that whenever you press \enquote{Compile} while working on a sub-chapter file, TeXstudio knows to compile the entire thesis starting from the main file.
		
		\subsection{Mathematical Assistance}
			For the complex equations required in engineering and technical theses, the `Mathematical Symbols' panel on the left provides a clickable library of Greek letters, operators, and relations. 
			Hovering over a symbol will even show you the required \LaTeX\ command.